% ==============================================================================
% OBJETIVOS
% ==============================================================================
% Contenido provisto directamente por el estudiante. El Objetivo específico
% 3 coincide verbatim con protocolo_sebastian_v1 (3).pdf, sección 4.1.
%
% Citas owasp2021 y mitreattack: el estudiante no adjuntó bibliografía esta
% vez. Se verificaron directamente contra las páginas oficiales por WebFetch
% antes de agregarlas al .bib (no se citaron de memoria):
%   - owasp2021: confirmado el título "A03:2021 -- Injection" y la URL
%     canónica en owasp.org.
%   - mitreattack: confirmado el nombre del framework y la organización
%     (MITRE); el sitio no expone un número de versión o fecha de
%     publicación único, se cita como framework de actualización continua.
%
% Corrección 2026-10-02: el Objetivo específico 1 afirmaba que los cuatro
% vectores caen bajo "OWASP Top 10 2021 (A03:2021)". Es incorrecto para
% Path Traversal/LFI, que está mapeado en A01:2021 -- Broken Access Control,
% no en A03:2021 -- Injection (verificado por WebFetch contra
% top10.owasp.org; ver corrección completa documentada en
% i-marco_teorico.tex). Se corrigió para distinguir ambas categorías en vez
% de agrupar los cuatro vectores bajo una sola.
\ifdefined\modoensamble
\else
	\input{0-preambulo_standalone}
\fi

\chapter{Objetivos}

\section{Objetivo general}
Diseñar e implementar una librería npm que integre un pipeline de ingeniería de características y
un modelo híbrido de aprendizaje automático capaz de detectar y clasificar amenazas conocidas y
anomalías estadísticas correspondientes a variantes desconocidas de los vectores cubiertos en
endpoints Node.js en tiempo de ejecución, con el propósito de proveer a startups tecnológicas sin
equipo de seguridad dedicado una herramienta accesible de protección a nivel de capa de
aplicación que no comprometa el rendimiento del servidor anfitrión.

\section{Objetivos específicos}
\begin{enumerate}
    \item Evaluar cuatro vectores de ataque intra-request dirigidos a endpoints Node.js — SQL
    Injection, XSS, Path Traversal/LFI y Command Injection — conforme a su prevalencia en el
    OWASP Top 10 2021 \cite{owasp2021} (SQL Injection, XSS y Command Injection bajo A03:2021 --
    Injection; Path Traversal/LFI bajo A01:2021 -- Broken Access Control) y su cobertura en la
    taxonomía MITRE ATT\&CK \cite{mitreattack}, para establecer las categorías de amenaza que
    estructurarán el dataset de entrenamiento del OE2 y los criterios de evaluación del OE3.

    \item Construir un entorno de prueba controlado y reproducible que genere un dataset
    estructurado con un mínimo de 100{,}000 muestras de tráfico HTTP etiquetadas; balanceando
    tráfico legítimo y ataques de las categorías seleccionadas, documentando los scripts de
    generación y la metodología de etiquetado, a fin de disponer de una base de datos de
    entrenamiento y validación que garantice la representatividad estadística necesaria para el
    entrenamiento de los modelos de detección.

    \item Evaluar la efectividad del sistema de detección implementado mediante precisión, recall
    y F1-score por categoría de ataque, verificando como criterio mínimo de aceptación un
    F1-score $\geq 0.80$ en al menos tres de las cuatro categorías seleccionadas bajo condiciones
    de tráfico simulado, umbral establecido en función del balance operacionalmente tolerable
    entre falsos negativos y falsos positivos en entornos de producción sin analista de seguridad
    dedicado.
\end{enumerate}

\ifdefined\modoensamble
\else
	\section*{Referencias}
	\printbibliography[heading=none]
	\end{document}
\fi
